% 007 业务模型文本：正文第 2 章输入。
\subsection{2.1 愿景}
\subsubsection{2.1.1 愿景与证据边界}
\BSRKeyPoint{组织愿景}{待 BSR 确认}{在第二门店及未来更多门店共用可信的客户、合同和单件库存记录，使顾客按预约时段快速取得合适、安全的装备，并让跨店归还和维修可追踪。}

\noindent 组织的关键涉众是经营负责人 Ole Fjordal 与度假村管理方。案例能证明第二租赁点扩张压力、缓解主店拥堵的意图和 Excel 单人操作局限；尚无 BSR 提供的排队、库存差错或维修周转基线。下表是建议指标，不能写成已完成绩效。

\begin{BSRTable}
\rowcolor{TableHead}\color{white}\bfseries 指标及口径 & \color{white}\bfseries 建议目标与状态 \\
预约可履约率：可用单件覆盖已确认预约的比例 & 建议 $\geq 98\%$；基线与目标待 BSR 确认。\\
到店取货时长：核身份至完成逐件交付的中位数 & 建议 $\leq 10$ 分钟；须现场测量。\\
库存账实一致率：盘点一致实物数/抽盘实物数 & 建议 $\geq 99\%$；须跨店抽盘验证。\\
疑似损坏初评及时率：发现后 24 小时内反馈的比例 & 建议 $\geq 95\%$；初评不是维修完成。\\
\end{BSRTable}

\subsubsection{2.1.2 执行者、工人与业务实体}
\begin{BSRTable}
\rowcolor{TableHead}\color{white}\bfseries 身份 & \color{white}\bfseries 价值、职责与证据边界 \\
业务执行者：顾客 & 取得合适装备、透明租金与押金、按约取还及可核对的异常费用。线上预约属目标蓝图。\\
内部关键涉众：经营负责人 Ole & 代表 BSR 经营目标，批准愿景、费率与跨店政策；其管理岗位在组织内部，不列为外部业务执行者。\\
外部关键涉众：度假村管理方 & 关注主店拥堵和第二租赁点扩张；现有证据未说明其向 BSR 请求何种具体业务服务，暂不连接业务用例。\\
业务工人：租赁专员 & 案例支持员工参与合同及故障登记；目标中负责备货、核身份/码、交付与逐件归还。\\
业务工人：门店经理 & 目标中负责排班、调拨、受损复核和安全验收；现行分工待访谈。\\
业务工人：装备检修技师/分包商 & 案例明确分包商定期到访并按纸单维修；在所建模的维修责任范围内作为业务工人，目标中按授权工单记录证据与工时。\\
业务实体：Excel 软件、目标 BSR 系统 & 表示能承担信息处理职责的非人系统；现状仅有 Excel 用途证据，BSR 系统属于改进蓝图。\\
业务资料与实物：客户记录、合同、装备、纸质维修单、货架 & 是处理对象或消息参数；不赋予纸张、装备或货架业务处理责任，不将它们放在业务序列生命线上。\\
\end{BSRTable}

\subsection{2.2 业务用例图}
\subsubsection{2.2.1 对外业务用例与价值}
\begin{BSRTable}
\rowcolor{TableHead}\color{white}\bfseries 业务用例 & \color{white}\bfseries 对执行者的可观察结果 \\
BU-01 租到合适、安全的装备 & 顾客在约定门店和时段取得适合尺码与水平的安全装备，知悉租金、押金及租赁责任。预约、报价、合同确认、付款与备货交付共同实现该价值。\\
BU-02 归还装备并取得结算 & 顾客得到逐件归还确认与可核对的租金、超时费、损坏费及押金结果；异常或争议有明确处理状态。异店归还属于该服务的约定方式。\\
\end{BSRTable}

\begin{figure}[htbp]\centering\resizebox{.94\linewidth}{!}{\BSRBusinessCases}
\caption{BU-01～02 BSR 对顾客提供的完整业务价值；内部维修与调拨属于实现支撑}\label{fig:b-cases}\end{figure}\FloatBarrier

\noindent B-CASE-01～05 保留为旧版实现片段索引，便于核对既有追踪材料，\textbf{不再表示五个独立的对外业务用例}。
\begin{BSRTable}
\rowcolor{TableHead}\color{white}\bfseries 原实现片段索引 & \color{white}\bfseries 现行业务价值映射 \\
B-CASE-01 预约与报价 & 实现 BU-01；对应 B-ASIS-01 / B-TOBE-01。\\
B-CASE-02 备货与交付 & 实现 BU-01；对应 B-ASIS-02 / B-TOBE-02。\\
B-CASE-03 接收归还 & 实现 BU-02；对应 B-ASIS-03 / B-TOBE-03。\\
B-CASE-04 维护装备 & 内部检修支撑 BU-01，受损结算证据支撑 BU-02；对应 B-ASIS-04 / B-TOBE-04。\\
B-CASE-05 跨店作业 & 异店归还实现 BU-02；永久调拨为保障 BU-01 的内部库存作业，对应 B-ASIS-05 / B-TOBE-05。\\
\end{BSRTable}

\subsubsection{2.2.2 图示规范与证据方法}
\noindent 本章依《软件方法：业务建模和需求（上，第 2 版）》印刷页 68～78 区分组织外的业务执行者、组织内承担职责的业务工人，以及外部执行者期望取得的完整价值。组织边界按责任划分，不能只凭工作地点、雇佣或外包形式判断；经营负责人和度假村的愿景利益另列为涉众。

\noindent 依该书印刷页 101～105，业务序列中 A 指向 B 的消息表示要求 B 承担职责；生命线的粒度是人或非人系统。纸合同、维修单和装备没有处理能力，作为消息参数表达；检查、搬运和实物验收由人员自调用表达。图标保留课程中的小人、圈中小人和空圆圈，分别表示业务执行者、业务工人与非人业务实体。

\noindent 实线实箭头表示职责请求，实线开箭头表示主动通知；水平虚线箭头只表示返回消息，垂直虚线表示生命线。分析图的暖色实线表示异常分支，\textbf{线型不表示证据可信度}。现状仅绘案例已支持的事实片段，证据不足时给出说明卡并在证据表列出待核问题；不能用虚线拼成一套未经观察的现行流程。改进图均为待 BSR 确认的目标协作，系统自动核验和提醒不替代人员的实物责任。
